% Options for packages loaded elsewhere
\PassOptionsToPackage{unicode}{hyperref}
\PassOptionsToPackage{hyphens}{url}
%
\documentclass[
  letterpaper,
]{scrbook}

\usepackage{amsmath,amssymb}
\usepackage{iftex}
\ifPDFTeX
  \usepackage[T1]{fontenc}
  \usepackage[utf8]{inputenc}
  \usepackage{textcomp} % provide euro and other symbols
\else % if luatex or xetex
  \usepackage{unicode-math}
  \defaultfontfeatures{Scale=MatchLowercase}
  \defaultfontfeatures[\rmfamily]{Ligatures=TeX,Scale=1}
\fi
\usepackage{lmodern}
\ifPDFTeX\else  
    % xetex/luatex font selection
\fi
% Use upquote if available, for straight quotes in verbatim environments
\IfFileExists{upquote.sty}{\usepackage{upquote}}{}
\IfFileExists{microtype.sty}{% use microtype if available
  \usepackage[]{microtype}
  \UseMicrotypeSet[protrusion]{basicmath} % disable protrusion for tt fonts
}{}
\makeatletter
\@ifundefined{KOMAClassName}{% if non-KOMA class
  \IfFileExists{parskip.sty}{%
    \usepackage{parskip}
  }{% else
    \setlength{\parindent}{0pt}
    \setlength{\parskip}{6pt plus 2pt minus 1pt}}
}{% if KOMA class
  \KOMAoptions{parskip=half}}
\makeatother
\usepackage{xcolor}
\ifLuaTeX
  \usepackage{luacolor}
  \usepackage[soul]{lua-ul}
\else
  \usepackage{soul}
  
\fi
\setlength{\emergencystretch}{3em} % prevent overfull lines
\setcounter{secnumdepth}{5}
% Make \paragraph and \subparagraph free-standing
\ifx\paragraph\undefined\else
  \let\oldparagraph\paragraph
  \renewcommand{\paragraph}[1]{\oldparagraph{#1}\mbox{}}
\fi
\ifx\subparagraph\undefined\else
  \let\oldsubparagraph\subparagraph
  \renewcommand{\subparagraph}[1]{\oldsubparagraph{#1}\mbox{}}
\fi


\providecommand{\tightlist}{%
  \setlength{\itemsep}{0pt}\setlength{\parskip}{0pt}}\usepackage{longtable,booktabs,array}
\usepackage{calc} % for calculating minipage widths
% Correct order of tables after \paragraph or \subparagraph
\usepackage{etoolbox}
\makeatletter
\patchcmd\longtable{\par}{\if@noskipsec\mbox{}\fi\par}{}{}
\makeatother
% Allow footnotes in longtable head/foot
\IfFileExists{footnotehyper.sty}{\usepackage{footnotehyper}}{\usepackage{footnote}}
\makesavenoteenv{longtable}
\usepackage{graphicx}
\makeatletter
\def\maxwidth{\ifdim\Gin@nat@width>\linewidth\linewidth\else\Gin@nat@width\fi}
\def\maxheight{\ifdim\Gin@nat@height>\textheight\textheight\else\Gin@nat@height\fi}
\makeatother
% Scale images if necessary, so that they will not overflow the page
% margins by default, and it is still possible to overwrite the defaults
% using explicit options in \includegraphics[width, height, ...]{}
\setkeys{Gin}{width=\maxwidth,height=\maxheight,keepaspectratio}
% Set default figure placement to htbp
\makeatletter
\def\fps@figure{htbp}
\makeatother

\makeatletter
\@ifpackageloaded{bookmark}{}{\usepackage{bookmark}}
\makeatother
\makeatletter
\@ifpackageloaded{caption}{}{\usepackage{caption}}
\AtBeginDocument{%
\ifdefined\contentsname
  \renewcommand*\contentsname{Table of contents}
\else
  \newcommand\contentsname{Table of contents}
\fi
\ifdefined\listfigurename
  \renewcommand*\listfigurename{List of Figures}
\else
  \newcommand\listfigurename{List of Figures}
\fi
\ifdefined\listtablename
  \renewcommand*\listtablename{List of Tables}
\else
  \newcommand\listtablename{List of Tables}
\fi
\ifdefined\figurename
  \renewcommand*\figurename{Figure}
\else
  \newcommand\figurename{Figure}
\fi
\ifdefined\tablename
  \renewcommand*\tablename{Table}
\else
  \newcommand\tablename{Table}
\fi
}
\@ifpackageloaded{float}{}{\usepackage{float}}
\floatstyle{ruled}
\@ifundefined{c@chapter}{\newfloat{codelisting}{h}{lop}}{\newfloat{codelisting}{h}{lop}[chapter]}
\floatname{codelisting}{Listing}
\newcommand*\listoflistings{\listof{codelisting}{List of Listings}}
\makeatother
\makeatletter
\makeatother
\makeatletter
\@ifpackageloaded{caption}{}{\usepackage{caption}}
\@ifpackageloaded{subcaption}{}{\usepackage{subcaption}}
\makeatother
\ifLuaTeX
  \usepackage{selnolig}  % disable illegal ligatures
\fi
\usepackage{bookmark}

\IfFileExists{xurl.sty}{\usepackage{xurl}}{} % add URL line breaks if available
\urlstyle{same} % disable monospaced font for URLs
\hypersetup{
  pdftitle={ThamesWater Andrio Emilio application},
  pdfauthor={Andrio Emilio},
  hidelinks,
  pdfcreator={LaTeX via pandoc}}

\title{ThamesWater Andrio Emilio application}
\author{Andrio Emilio}
\date{}

\begin{document}
\frontmatter
\maketitle

\renewcommand*\contentsname{Table of contents}
{
\setcounter{tocdepth}{2}
\tableofcontents
}
\mainmatter
\bookmarksetup{startatroot}

\chapter*{Preface}\label{preface}
\addcontentsline{toc}{chapter}{Preface}

\markboth{Preface}{Preface}

Hi!

I'm Andrio, I received a pre-call screening for the performance and risk
analyst role here at ThamesWater. I wanted to be proactive and have all
my research, preparation and skills in this Quarto book. This is meant
to be candid and for my personal use. As such it will show my thought
process between each chapter and also my (amazing) personality.\\
\strut \\
I hope whoever reads this enjoys (including myself)!

:)

\bookmarksetup{startatroot}

\chapter{Past experience}\label{past-experience}

On this page I wanted to briefly map out some past experience and skills
in accordance to the job specification.

I've stuck on a candle and a famous Andrio Emilio spotify playlist so
lets go!

\section{Analytical Background}\label{analytical-background}

\subsection{Excel}\label{excel}

I have been using excel in the workplace for nearly 10 years now. In my
first role I was a data researcher and processor. This involved
collecting ONS \& local authority data, transforming, cleaning and
ultimately providing analysis with visualisations. I primarily used
vlookup to join time series data with the next entry. Conditional
formatting to track differences and bespoke analysis with pivot charts
and pivot tables. This gave me the foundations to transition towards my
first analyst role at imperial college London. A hands on role to which
I had to give more application and thought towards the processes
required to transform and present highly sensitive HR data. Alongside
this, I was tasked with more and more data to handle and had to learn
automation skills. This led me to developing macros for the use of
initial data transformation and soon introduced me to powerBI. Now I use
excel sparingly, I have tracker templates to record my response rates
for tasks as well as a personal diary. I have also recoreded a
systematic review in excel (a database recording features from over 100
academic papers) .

\subsection{SQL}\label{sql}

I have used SQL in tandem with power BI to connect and refresh data for
the use of weekly, monthly and annual reports at my role in Network
rail. I have completed several modules during my masters study that used
SQL. These focused upon creating spatial databases as well as spatial
joins. Since, I have completed a 7 week analyst course in which
practical application of SQL was used. This involved using SQL to
connect to databases, create data, create tables, transfor data, merge,
join \& append.

\subsection{Power BI}\label{power-bi}

I have numerous years experience in power BI. This started in my role as
a hr information insight analyst at imperial college. In this I created
bespoke reports based upon foi requests as well as monthly HR necessary
data (such as PTO, sickness and diversity metrics). In my role at
network rail I used power BI top create asset based metrics for
stakeholders decision making. An example of a project from start to
finish would be a vegetation risk asset PBI model report. This tracked
from live recorded mothly input data of vegetation risk on the railway.
It covered species, protected type and seasonality. It also recorded
previous intervention \& planned removal. These were refreshed with
every new report (usually an excel document) that would then be
transformed in M-query to handle anomalous values, different data types
and wrong inputs. For stakeholders, they would a live map showing
new/planned vegetation management and metric charts \& tables to show if
they are meeting required quotas.

\begin{itemize}
\tightlist
\item
  Asset management project you did at network rail to assess performance
  of track. Go in depth with charts, spatial visualization and
  automation and how this benefited stakeholders across the company.
\end{itemize}

\section{Other technical experience}\label{other-technical-experience}

\subsection{Python}\label{python}

I have used Python for model creation and analysis of data in my masters
studies at UCL. My dissertation project was computed in python and this
involved the use of unsupervised/supervised models such as random forest
and k means clusters. I have also done some regression analysis and used
libraries such as pandas, numpy and scikit learn. Because of this I am
used to working in jupyter notebooks and would love the opportunity to
pursue this further

\subsection{Geospatial tools}\label{geospatial-tools}

My geosptaial knowledge started within my role at network rail and led
me to pursue a masters at UCL. I completed an Msc in Urban spatial
science in which gave me a strong foundation in geospatial analysis and
also allowed me to improve my autonomous skills to meet deadline and
choose suitable projects. Because of this I have used a vast range of
geospatial data and have a knowledge of many areas, including networks,
remote sensing and data science.

I used geospatial tools in an applied asset-heavy environment at
\textbf{Network Rail}. There, I worked with GIS specialists on mapping
over 3000 miles of track for the Southern route, helping turn asset and
engineering data into more useful visual reporting. I used tools such as
\textbf{ArcGIS, QGIS}, and map-based visualisation approaches, and I
also integrated \textbf{Mapbox visuals into Power BI reports} so
stakeholders could see issues spatially rather than just in tables. That
was especially useful for visualising track-related issues and helping
teams understand where problems or risk-related areas were concentrated
across the network.

\subsection{MS access \& databricks}\label{ms-access-databricks}

\begin{itemize}
\item
  no formal application but have planned time on courses starting next
  week.

  Training for access:
  \href{https://support.microsoft.com/en-us/office/access-video-training-a5ffb1ef-4cc4-4d79-a862-e2dda6ef38e6}{here}

  Training for databricks:
  \href{https://learn.microsoft.com/en-us/training/paths/data-engineer-azure-databricks/}{here}
\end{itemize}

\section{Asset based environments}\label{asset-based-environments}

\subsection{BI reports}\label{bi-reports}

I\textquotesingle ve worked in an asset-based environment through my
role at \textbf{Network Rail}, where the data and reporting were closely
tied to physical infrastructure and operational performance. That meant
working with engineering and operational datasets linked to the rail
network, and helping teams understand where issues, risks, or
constraints were appearing across the assets. I think that experience is
relevant because it taught me how to use data not just for reporting,
but to support decisions in a complex infrastructure setting

One of my main responsibilities at Network Rail was supporting Kent and
Sussex track teams through BI reporting. I integrated over 15 data
sources into Power BI and used performance measures such as KPIs, speed
restrictions, and risk-related reporting to help stakeholders understand
how the network was performing. Stakeholders could move away from manual
reporting and instead use a more consistent view of performance to
prioritise work and track issues over time.

\section{Maintenance \& communication}\label{maintenance-communication}

\begin{itemize}
\item
  Maintaining existing data management processes, how have you done this
  before?

  \begin{quote}
  Making sure I understand the process end to end, where data comes
  from, how its transformed, who will be using and where are the risks.

  You have had a Microsoft one note diary recording of your systematic
  review to track when changes were made and what the change was.

  Understanding of existing data management processes by communicating
  and organising meetings with data teams and stakeholders to understand
  unfamiliar data

  Creating very transparent and available documents for stakeholders and
  personal use (available in sharepoint) to explain data

  Monitoring key data quality metrics in automated BI report that
  assesses weekly, quarterly and yearly performance against previous
  entry.
  \end{quote}
\item
  What metrics did you use to present this information? List all the
  tools, charts and visualisations.

  \begin{quote}
  \subsection{Data quality metrics}\label{data-quality-metrics}

  \begin{itemize}
  \item
    missing values / null counts
  \item
    duplicate records
  \item
    invalid or out-of-range values
  \item
    consistency of categories / naming
  \item
    unexpected period changes
  \item
    completeness by source or region
  \item
    refresh success / failed updates
  \item
    record counts compared with previous periods
  \item
    percentage change against prior week / quarter / year
  \item
    exception counts / flagged anomalies
  \end{itemize}

  \subsection{Performance metrics}\label{performance-metrics}

  \begin{itemize}
  \item
    KPI movement over time
  \item
    period-on-period comparisons
  \item
    count of issues / restrictions / events
  \item
    trend direction
  \item
    thresholds breached
  \item
    concentration by area / route / asset group
  \item
    backlog / open vs resolved items
  \item
    performance against target
  \end{itemize}
  \end{quote}
\item
  List example of collaboration

  \begin{quote}
  Network rail vegetation BI report, how you used an iterative
  continuous improvement process
  \end{quote}
\end{itemize}

\section{Producing detailed reports}\label{producing-detailed-reports}

I\textquotesingle ve produced detailed reports in both professional and
academic settings. I\textquotesingle m comfortable adapting the level of
detail depending on the audience. At Network Rail, that involved
developing reports and presentations based on engineering and
operational data. As the aim was to make complex information clear
enough for stakeholders to act on. It was important to present the
detail accurately while still making the key messages easy to
understand.

During my master\textquotesingle s, I also worked on group projects
where we had to take technical analysis and present it in a structured
and professional way. That helped me build confidence in explaining
analytical work clearly. In my dissertation, I worked closely with
industry experts to shape the focus of the project and make sure the
outputs were relevant as well as technically robust. That experience
taught me that producing a strong report is not just about analysis
it\textquotesingle s also about understanding the audience,
communicating clearly, and making the findings useful.

\bookmarksetup{startatroot}

\chapter{Research}\label{research}

An exciting part of this process is finding out information I didn't
already know about Thames Water. I wanted to understand the company
through the Andrio Emilio lense and find someinteresting case studies!

\section{Company overview}\label{company-overview}

\emph{``Every day, we serve 16 million customers across London and the
Thames Valley. From that first cup of tea in the morning through to
baths at bedtime.}

\emph{As well as supplying water to your taps, we take wastewater away.
We collect and treat sewage, returning clean water to our rivers.}

\emph{Delivering~\href{https://www.thameswater.co.uk/about-us/thames-water-explained}{life\textquotesingle s
essential service}, so our customers, communities and environment can
thrive. We're here to keep taps flowing for generations to come. Find
out more about how
we're~\href{https://www.thameswater.co.uk/about-us/our-biggest-upgrade}{upgrading
our network}.''}

It is a very informative site, compiling pollution reduction,
environmental information, five year plans and assessing performance.

\begin{itemize}
\item
  Open data strategy, so data is freely available for everyone to
  access.
\item
  PR24 five year plan 2025 - 2030. Involving strategy, cost efficiency,
  expected outcomE and data models.
\item
  Biggest upgrade in 150 years with 20 planned investment focused on
  reducing pollution protecting water quality and installing over 1
  million more smart meters.
\item
  Digital pages talk about building an \textbf{enterprise data platform}
  on \textbf{Databricks Lakehouse}
\end{itemize}

\section{Team focus}\label{team-focus}

Why is an asset performance analyst essential?

\emph{``Asset performance metrics form the foundation of data-driven
maintenance, providing continuous visibility into equipment reliability,
maintenance efficiency, and operational trends that would otherwise
remain hidden.}

\emph{Tracking the right metrics connects maintenance activities to
business outcomes, helping teams reduce unplanned downtime, control
costs, extend equipment life, and demonstrate measurable value to
leadership.}

\emph{Advanced software platforms automate metric collection and
analysis, replacing manual spreadsheets with real-time dashboards,
AI-powered insights, and integrated workflows that turn data into
action.}''

\textbf{Assessing performance and behaviour, what does this mean for
Thames Water?}

\ul{Performance is about how well the asset is doing}

\ul{Behaviour is about the pattern of how the asset acts\\
}

For performance:

\begin{itemize}
\item
  how well is it doing?
\item
  is it meeting expected standards?
\item
  how many incidents / failures / downtime events are there?
\end{itemize}

For behavioural:

\begin{itemize}
\item
  what pattern does it show?
\item
  when does it tend to go wrong?
\item
  under what conditions does it behave differently?
\item
  is there a recurring operating signature that suggests risk?
\end{itemize}

\subsection{Why this is important}\label{why-this-is-important}

\section{Case studies}\label{case-studies}

\section{Policy, ethics and
competences}\label{policy-ethics-and-competences}

\begin{itemize}
\tightlist
\item
  \textbf{Take care}
\end{itemize}

\begin{itemize}
\tightlist
\item
  \textbf{Passionate about everything we do}
\end{itemize}

\begin{itemize}
\tightlist
\item
  \textbf{Be respectful and value everyone}
\end{itemize}

\begin{itemize}
\tightlist
\item
  \textbf{Reach higher, be better}
\end{itemize}

\begin{itemize}
\tightlist
\item
  \textbf{Take ownership}
\end{itemize}

\begin{itemize}
\tightlist
\item
  \textbf{Be proud, Be blue / Team Thames}
\end{itemize}

\bookmarksetup{startatroot}

\chapter{Preparation}\label{preparation}

\section{Questions}\label{questions}

\subsection{1 "Can you talk me through your
background?"}\label{can-you-talk-me-through-your-background}

I started working in data during my undergraduate degree. In my first
year I scored 100\% in a business statistics module and this led me to
change my degree to analytics. I then sought experience in a Leeds based
startup alongside my degree during the second and 3rd year. In this role
I was a data researcher and processor. This meant collecting, cleaning,
transforming and creating charts for councils using ONS and local
authority data. Because of this experience in excel I started my first
graduate role and imperial college london. Wherein in the first week my
manager gave me a huge power BI book, said learn as much as you can and
by the end of the second week I want to see a report from you. This led
me to learn power BI and has been a staple for my work. It was
especially important as i started my role at network rail, learning
asset and engineering data, how to visualise these metrics and
understand what data insights stakeholders want to see! With this I got
the oportunity to work alongside a lead geospatial scientist, and led to
a power BI project of mapping asset faults spatially. This introduced me
to geospatial analysis in which I studied a masters in at UCL and scored
a distinction.

\subsection{2 "Why are you interested in this
role?"}\label{why-are-you-interested-in-this-role}

Im interested in the role as it has both familiarity in working within
an asset data based environment but also the opportunity to learn how to
apply the skills I have learnt and be a part of a winder data team to
work for real world data problems to which I can see a clear benefit. I
think in the past, especially during my masters degree, I learnt that
working alongside stakeholders and experts who wanted to know more about
the data in an industry they have expertise in knowing I could do just
that. I also think this is a position where I can develop my analytical
skills further and keep learning, especially within databricks.

I\textquotesingle m also drawn to the scale and complexity of the
organisation. It\textquotesingle s the kind of environment where you
have to think carefully, work collaboratively, and make decisions that
are both practical and responsible. That really appeals to me, because I
enjoy roles where I can contribute in a structured way, take ownership
of my work, and help improve how things are done.

\subsection{3 "What attracted you to a role in asset operations /
wastewater /
utilities?"}\label{what-attracted-you-to-a-role-in-asset-operations-wastewater-utilities}

The basis for my application is my prior skills in an asset based
environment at network rail. Since this role I have explored other areas
within data, including geospatial analysis and research intensive roles.
Now I feel that my skills in insight, visualisation and communicating
this is well geared toward a role at Thames Water

\subsection{4 "This is a 1-year fixed-term contract are you comfortable
with
that?}\label{this-is-a-1-year-fixed-term-contract-are-you-comfortable-with-that}

Yes, this gives me a goal to garner and hone my skills before i commit
to the next opportunity.

\subsection{5 "How strong are you in
SQL?"}\label{how-strong-are-you-in-sql}

I have experience, creating, querying and transforming data in SQL. In
the past I have used SQL to spatially join, extract data for power BI
usage and experiment. I have developed a small project on this page to
show my basic understanding of SQL. I have also undergone refresher
courses to make sure I have the baseline for success in this role.

\subsection{\texorpdfstring{6 \textbf{"What kinds of queries /
transformations have you
written?"}}{6 "What kinds of queries / transformations have you written?"}}\label{what-kinds-of-queries-transformations-have-you-written}

\subsection{7 "How have you used Power BI in previous
roles?"}\label{how-have-you-used-power-bi-in-previous-roles}

STAR Example

Situation

Task

Action

Result

\subsection{8 "What dashboards or management reports have you
built?"}\label{what-dashboards-or-management-reports-have-you-built}

S

T

A

R

\subsection{9 "Tell me about a time you turned complex analysis into an
actionable
recommendation."}\label{tell-me-about-a-time-you-turned-complex-analysis-into-an-actionable-recommendation.}

S

T

A

R

\subsection{10 "How do you present data to non-technical
stakeholders?"}\label{how-do-you-present-data-to-non-technical-stakeholders}

S

T

A

R

\subsection{11 "How do you handle situations where stakeholders want
different things from the
data?"}\label{how-do-you-handle-situations-where-stakeholders-want-different-things-from-the-data}

S

T

A

R

\subsection{12 "Can you give an example of influencing action through
reporting or
insight?"}\label{can-you-give-an-example-of-influencing-action-through-reporting-or-insight}

\subsection{13 How do you deal with data quality
problems?}\label{how-do-you-deal-with-data-quality-problems}

\subsection{14 Have you worked with operational / asset / infrastructure
data
before?}\label{have-you-worked-with-operational-asset-infrastructure-data-before}

\subsection{15 Tell me about a time you improved data quality or a
reporting
process.}\label{tell-me-about-a-time-you-improved-data-quality-or-a-reporting-process.}

\subsection{16 Tell me about a time you identified a risk or performance
issue from
data.}\label{tell-me-about-a-time-you-identified-a-risk-or-performance-issue-from-data.}

\subsection{17 "What kinds of statistical modelling have you
done?"}\label{what-kinds-of-statistical-modelling-have-you-done}

My experience ranges from practical operational analysis to more formal
modelling. In academic and research settings, I\textquotesingle ve used
Python-based modelling, geospatial analysis and machine learning
approaches to identify patterns and relationships in data. In more
operational roles, statistical modelling has been more about trend
analysis, KPI monitoring, anomaly detection, and identifying where
performance patterns may indicate risk or require further investigation.

\subsection{\texorpdfstring{18 \textbf{The role mentions statistical
modelling. What does that mean to you in
practice?"}}{18 The role mentions statistical modelling. What does that mean to you in practice?"}}\label{the-role-mentions-statistical-modelling.-what-does-that-mean-to-you-in-practice}

\textbf{Response:}\\
"To me, statistical modelling in practice is about using data in a
structured way to understand performance, identify patterns, and support
decisions. It can include more advanced predictive techniques, but in
many operational roles it also means applying methods like trend
analysis, KPI monitoring, and anomaly detection. Trend analysis is about
looking at how a metric changes over time to understand whether
performance is stable, improving, or deteriorating. KPI monitoring is
about tracking the most important measures against targets, thresholds,
or expected ranges so it is clear where performance is on track and
where it may need attention. Anomaly detection is about spotting values
or patterns that fall outside what would normally be expected, which can
help flag possible operational issues, emerging risk, or data quality
problems. So overall, I\textquotesingle d describe statistical modelling
as using data systematically to move beyond simple reporting and help
identify where further investigation or action may be needed."

\subsection{\texorpdfstring{\textbf{19 Can you explain a time when you
had to analyze a complex set of data and make recommendations based on
your
findings?}}{19 Can you explain a time when you had to analyze a complex set of data and make recommendations based on your findings?}}\label{can-you-explain-a-time-when-you-had-to-analyze-a-complex-set-of-data-and-make-recommendations-based-on-your-findings}

\subsection{\texorpdfstring{\textbf{20. How would you handle
disagreements with management about your
analysis?}}{20. How would you handle disagreements with management about your analysis?}}\label{how-would-you-handle-disagreements-with-management-about-your-analysis}

\subsection{\texorpdfstring{\textbf{21 What methods do you use to ensure
data
accuracy?}}{21 What methods do you use to ensure data accuracy?}}\label{what-methods-do-you-use-to-ensure-data-accuracy}

\subsection{\texorpdfstring{\textbf{22 How do you ensure that your
analytics reports are easily understood by non-analytic
professionals?}}{22 How do you ensure that your analytics reports are easily understood by non-analytic professionals?}}\label{how-do-you-ensure-that-your-analytics-reports-are-easily-understood-by-non-analytic-professionals}

\subsection{\texorpdfstring{\textbf{23 How do you ensure that your
analytics reports are easily understood by non-analytic
professionals?}}{23 How do you ensure that your analytics reports are easily understood by non-analytic professionals?}}\label{how-do-you-ensure-that-your-analytics-reports-are-easily-understood-by-non-analytic-professionals-1}

\subsection{\texorpdfstring{\textbf{24 What experience do you have with
data visualization
tools?}}{24 What experience do you have with data visualization tools?}}\label{what-experience-do-you-have-with-data-visualization-tools}

\subsection{\texorpdfstring{\textbf{25 How do you handle large data
sets?}}{25 How do you handle large data sets?}}\label{how-do-you-handle-large-data-sets}

\subsection{\texorpdfstring{\textbf{26 Can you describe a situation
where you had to use analytical skills to solve a
problem?}}{26 Can you describe a situation where you had to use analytical skills to solve a problem?}}\label{can-you-describe-a-situation-where-you-had-to-use-analytical-skills-to-solve-a-problem}

\subsection{\texorpdfstring{\textbf{27 Can you explain a time when your
analysis led to a considerable improvement in business
performance?}}{27 Can you explain a time when your analysis led to a considerable improvement in business performance?}}\label{can-you-explain-a-time-when-your-analysis-led-to-a-considerable-improvement-in-business-performance}

\subsection{\texorpdfstring{\textbf{28 What techniques do you use in
qualitative data
analysis?}}{28 What techniques do you use in qualitative data analysis?}}\label{what-techniques-do-you-use-in-qualitative-data-analysis}

\subsection{\texorpdfstring{\textbf{29 How do you prioritize your
analysis when dealing with multiple
projects?}}{29 How do you prioritize your analysis when dealing with multiple projects?}}\label{how-do-you-prioritize-your-analysis-when-dealing-with-multiple-projects}

\subsection{\texorpdfstring{\textbf{30 How comfortable are you with
presenting your findings to a senior management
team?}}{30 How comfortable are you with presenting your findings to a senior management team?}}\label{how-comfortable-are-you-with-presenting-your-findings-to-a-senior-management-team}

\subsection{\texorpdfstring{\textbf{31 How do you validate the
reliability and consistency of your
data?}}{31 How do you validate the reliability and consistency of your data?}}\label{how-do-you-validate-the-reliability-and-consistency-of-your-data}

Tell us about a time you did team work? Tell us when you used
initiative? Tell us about a time you had to draft something from
scratch?

\subsection{\texorpdfstring{\textbf{32 How have you monitored data
quality metrics in
reporting?"}}{32 How have you monitored data quality metrics in reporting?"}}\label{how-have-you-monitored-data-quality-metrics-in-reporting}

\textbf{Response:}\\
"I\textquotesingle ve approached that by building reporting that does
not only show performance outcomes, but also helps assess whether the
data and metrics are behaving as expected over time. For example,
I\textquotesingle d look at consistency across weekly, quarterly, and
yearly reporting periods, compare current entries against previous ones,
and flag unexpected changes, missing values, unusual spikes, or
inconsistencies. The purpose is to make it easier to distinguish between
a genuine operational change and something that may be a data issue."

\subsection{\texorpdfstring{\textbf{33 What have you done to keep data
processes organised and
reliable?"}}{33 What have you done to keep data processes organised and reliable?"}}\label{what-have-you-done-to-keep-data-processes-organised-and-reliable}

\textbf{Response:}\\
"I usually approach that in three parts. First, I build understanding by
speaking to the people closest to the data, especially when a dataset or
process is unfamiliar. That helps me understand the source, logic, and
any known issues. Second, I document things clearly, both for myself and
for others, so there\textquotesingle s transparency around definitions,
changes, and assumptions. For example, I\textquotesingle ve kept
structured change records in OneNote and created clear reference
documents for stakeholders in shared spaces like SharePoint. Third, I
monitor the outputs through reporting and data quality checks, so if
something changes or degrades, it can be picked up early."

\bookmarksetup{startatroot}

\chapter{EXCEL}\label{excel-1}

\begin{itemize}
\item
  Undertaking performance and behavioural analysis on waste network
  assets
\item
  ~Interventions are actively tracked to ensure timely and effective
  delivery against planned outcomes.
\item
  Using thameswater total daily flow data - multiple years.

  \begin{itemize}
  \item
    How would you clean this data set, think by column, transforming the
    data types making sure is numerical etc etc, making sure data is
    input correctly.
  \item
    How would you then make use of a vlookup
  \item
    Make use of conditional formatting
  \item
    Make use of pivot tables and charts
  \item
    Then introduce multiple years.
  \end{itemize}
\end{itemize}

\bookmarksetup{startatroot}

\chapter{SQL}\label{sql-1}

\bookmarksetup{startatroot}

\chapter{Power BI}\label{power-bi-1}

\begin{itemize}
\item
\item
  Produce a small powerBI report that is fully automated, self refreshes
  and has rules to send to stakeholders. Show how you connect to SQL and
  excel and how you use M-query to transform data before using DAX. Set
  text rules to provide contextual information of data ie `this week \_
  location has under performed in comparison to all other locations'
\item
  interpreting asset and operational data to identify performance trends
  and risk insights, and sharing findings using clear visualisations and
  commentary
\end{itemize}


\backmatter

\end{document}
